% part: intuitionistic-logic
% chapter: soundness-completeness
% section: soundness-nd

\documentclass[../../../include/open-logic-section]{subfiles}

\begin{document}

\olfileid{int}{sc}{snd}

\olsection{नैसर्गिक निगमनाची निर्दोषता}

आता संबंधी चिन्हार्थमीमांसेच्या दृष्टीने नैसर्गिक निगमनाची
निर्दोषता सिद्ध करू. म्हणजेच, गृहीतकांच्या संचापासून
एखादे !!a{formula} !!{derivable} असेल, तर त्या
गृहीतकांच्या संचापासून ते !!{formula}
तार्किकरीत्या निष्पन्न होते, हे दाखवू.

\begin{thm}[निर्दोषता]
  \ollabel{thm:soundness}
  $\Gamma \Proves !A$ असेल, तर $ \Gamma \Entails !A$.
\end{thm}

\begin{proof}
  $\Gamma \Proves !A$ असेल, तर $\Gamma \Entails !A$,
  हे सिद्ध करू. $\Gamma$ पासून~$!A$ च्या
  !!{derivation} वर आगमनाने सिद्धता करतो.

  \begin{enumerate}
  \item !!{derivation} मध्ये केवळ गृहीतक~$!A$ असेल,
    तर $!A \Proves !A$. आता $!A \Entails !A$ हे
    दाखवायचे आहे. $ \mSat{M}{!A}[w]$ असे समजा.
    मग उघडच $\mSat{M}{!A}[w]$.

  \item !!{derivation} चा शेवट $\Intro{\land}$ ने होतो:
    \iftag{probAnd}{सराव.}{आधारविधान~$!B$ ची
      !!{undischarged} गृहीतके~$\Gamma$ यांपासून आणि
      आधारविधान~$!C$ ची !!{undischarged}
      गृहीतके~$\Delta$ यांपासून असलेल्या
      !!{derivation}s वरून $\Gamma \Proves
      !B$ आणि $\Delta \Proves !C$ हे मिळते.
      आगमन-गृहीतकानुसार $\Gamma \Entails !B$
      आणि $\Delta \Entails !C$. संपूर्ण
      !!{derivation} ची !!{undischarged} गृहीतके
      $\Gamma$ व $\Delta$ मिळून असल्यामुळे
      $\Gamma \cup \Delta \Entails !B \land !C$
      हे दाखवायचे आहे. म्हणून
      $\mSat{M}{\Gamma \cup \Delta}[w]$ असे समजा.
      मग $\mSat{M}{\Gamma}[w]$ सुद्धा.
      $\Gamma \Entails !B$ असल्याने
      $\mSat{M}{!B}[w]$. त्याचप्रमाणे
      $\mSat{M}{!C}[w]$. म्हणून
      $\mSat{M}{!B \land !C}[w]$.}

  \item !!{derivation} चा शेवट $\Elim{\land}$ ने होतो:
    \iftag{probAnd}{सराव.}{आधारविधान $!B
        \land !C$ ची !!{undischarged} गृहीतके
        $\Gamma$ यांपासून असलेली !!{derivation}
        $\Gamma \Proves !B \land !C$ हे दाखवते.
        आगमन-गृहीतकानुसार
        $\Gamma \Entails !B \land !C$.
        $\Gamma \Entails !B$ हे दाखवायचे आहे.
        म्हणून $\mSat{M}{\Gamma}[w]$ असे समजा.
        $\Gamma \Entails !B \land !C$
        असल्याने $\mSat{M}{!B \land !C}[w]$.
        मग $\mSat{M}{!B}[w]$ सुद्धा.
        त्याचप्रमाणे $\Elim\land$ चा शेवट~$!C$
        ने होत असेल, तर $\Gamma \Entails !C$.}

  \item !!{derivation} चा शेवट $\Intro{\lor}$ ने होतो:
    \iftag{probOr}{सराव.}{आधारविधान $!B$ आहे,
      आणि~$!B$ येथे संपणाऱ्या !!{derivation} ची
      !!{undischarged} गृहीतके $\Gamma$ आहेत,
      असे समजा. मग $\Gamma \Proves !B$;
      आगमन-गृहीतकानुसार $\Gamma \Entails !B$.
      $\Gamma \Entails !B \lor !C$ हे दाखवायचे आहे.
      $\mSat{M}{\Gamma}[w]$ असे समजा.
      $\Gamma \Entails !B$ असल्याने
      $\mSat{M}{!B}[w]$. पण मग
      $\mSat{M}{!B \lor !C}[w]$ सुद्धा.
      त्याचप्रमाणे आधारविधान~$!C$ असेल, तर
      $\Gamma \Entails !C$.}

  \item !!{derivation} चा शेवट~$\Elim{\lor}$ ने होतो:
    \iftag{probOr}{सराव.}{आधारविधानांवर संपणाऱ्या
      !!{derivation}s अनुक्रमे $!B \lor !C$ ची
      !!{undischarged} गृहीतके~$\Gamma$,
      $!D$ ची !!{undischarged} गृहीतके
      $\Delta_1 \cup \{!B\}$, आणि $!D$ ची
      !!{undischarged} गृहीतके
      $\Delta_2 \cup \{!C\}$ यांपासून आहेत.
      म्हणून $\Gamma \Proves !B \lor !C$,
      $\Delta_1 \cup \{!B\} \Proves !D$,
      आणि $\Delta_2 \cup \{!C\} \Proves !D$.
      आगमन-गृहीतकानुसार
      $\Gamma \Entails !B \lor !C$,
      $\Delta_1 \cup \{!B\} \Entails !D$,
      आणि $\Delta_2 \cup \{!C\} \Entails !D$.
      $\Gamma \cup \Delta_1 \cup \Delta_2 \Entails !D$
      हे सिद्ध करायचे आहे.

    $\mSat{M}{\Gamma \cup \Delta_1 \cup \Delta_2}[w]$
    असे समजा. मग $\mSat{M}{\Gamma}[w]$;
    $\Gamma \Entails !B \lor !C$ असल्याने
    $\mSat{M}{!B \lor !C}[w]$.
    $\mSat{M}{}$ च्या व्याख्येनुसार
    $\mSat{M}{!B}[w]$ किंवा
    $\mSat{M}{!C}[w]$. म्हणून दोन प्रकार पाहू:
    (a)~$\mSat{M}{!B}[w]$. मग
    $\mSat{M}{\Delta_1 \cup \{!B\}}[w]$.
    $\Delta_1 \cup \{!B\} \Entails !D$
    असल्याने $\mSat{M}{!D}[w]$.
    (b)~$\mSat{M}{!C}[w]$. मग
    $\mSat{M}{\Delta_2 \cup \{!C\}}[w]$.
    $\Delta_2 \cup \{!C\} \Entails !D$
    असल्याने $\mSat{M}{!D}[w]$.
    म्हणून दोन्ही प्रकारांत
    $\mSat{M}{!D}[w]$, हेच दाखवायचे होते.}

  \item !!{derivation} चा शेवट $\Intro{\lif}$ ने होतो
    आणि निष्कर्ष~$!B\lif !C$ आहे. मग
    आधारविधान~$!C$ आहे, आणि त्या आधारविधानावर
    संपणाऱ्या !!{derivation} ची
    !!{undischarged} गृहीतके~$\Gamma \cup
    \{!B\}$ आहेत. म्हणून
    $\Gamma \cup \{!B\} \Proves !C$;
    आगमन-गृहीतकानुसार
    $\Gamma \cup \{!B\} \Entails !C$.
    $\Gamma \Entails !B \lif !C$ हे दाखवायचे आहे.

    $\mSat{M}{\Gamma}[w]$ असे समजा.
    $Rww'$ असे असलेल्या सर्व $w'$ साठी
    $\mSat{M}{!B}[w']$ असेल, तर
    $\mSat{M}{!C}[w']$, हे दाखवायचे आहे.
    म्हणून $Rww'$ आणि $\mSat{M}{!B}[w']$
    असे गृहीत धरा.
    \olref[sem][rel]{prop:true-monotonic} नुसार
    $\mSat{M}{\Gamma}[w']$.
    $\Gamma \cup \{!B\} \Entails !C$
    असल्याने $\mSat{M}{!C}[w']$,
    हेच दाखवायचे होते.

  \item !!{derivation} चा शेवट $\Elim{\lif}$ ने होतो
    आणि निष्कर्ष~$!C$ आहे. आधारविधाने
    $!B \lif !C$ आणि $!B$ आहेत;
    त्यांची !!{derivation}s अनुक्रमे
    !!{undischarged} गृहीतके $\Gamma$ आणि
    $\Delta$ यांपासून आहेत. म्हणून
    $\Gamma \Proves !B \lif !C$ आणि
    $\Delta \Proves !B$.
    आगमन-गृहीतकानुसार
    $\Gamma \Entails !B \lif !C$
    आणि $\Delta \Entails !B$.
    $\Gamma \cup \Delta \Entails !C$
    हे दाखवायचे आहे.

    $\mSat{M}{\Gamma \cup \Delta}[w]$ असे समजा.
    $\mSat{M}{\Gamma}[w]$ आणि
    $\Gamma \Entails !B \lif !C$ असल्याने
    $\mSat{M}{!B \lif !C}[w]$.
    व्याख्येनुसार, $Rww'$ असे असलेल्या
    सर्व $w'$ साठी $\mSat{M}{!B}[w']$
    असेल, तर $\mSat{M}{!C}[w']$.
    $R$ परावर्ती असल्याने $w$ हे
    $Rww'$ असलेल्या $w'$ पैकी एक आहे.
    म्हणजे $\mSat{M}{!B}[w]$ असेल,
    तर $\mSat{M}{!C}[w]$.
    $\mSat{M}{\Delta}[w]$ आणि
    $\Delta \Entails !B$ असल्याने
    $\mSat{M}{!B}[w]$. म्हणून
    अपेक्षेप्रमाणे $\mSat{M}{!C}[w]$.

  \item !!{derivation} चा शेवट $\FalseInt$ ने होतो
    आणि निष्कर्ष~$!A$ आहे. आधारविधान
    $\lfalse$ आहे; त्या आधारविधानाची
    !!{derivation} ज्या
    !!{undischarged} गृहीतकांपासून आहे,
    ती~$\Gamma$ आहेत. मग
    $\Gamma \Proves \lfalse$.
    आगमन-गृहीतकानुसार
    $\Gamma \Entails \lfalse$.
    $\Gamma \Entails !A$ हे दाखवायचे आहे.

    अप्रत्यक्ष सिद्धता करू.
    $\Gamma \Entails/ !A$ असेल, तर
    असे प्रतिमान~$\mModel{M}$ आणि जग~$w$
    आहेत की $\mSat{M}{\Gamma}[w]$
    आणि $\mSat/{M}{!A}[w]$.
    $\Gamma \Entails \lfalse$
    असल्याने $\mSat{M}{\lfalse}[w]$.
    पण हे अशक्य आहे; कारण व्याख्येनुसार
    $\mSat/{M}{\lfalse}[w]$.
    म्हणून $\Gamma \Entails !A$.
  \item !!{derivation} चा शेवट $\Intro\lnot$ ने होतो: सराव.
  \item !!{derivation} चा शेवट $\Elim\lnot$ ने होतो: सराव.
  \end{enumerate}
\end{proof}

\begin{prob}
  \olref[int][sc][snd]{thm:soundness} ची सिद्धता
  पूर्ण करा. $\Intro\lnot$ आणि $\Elim\lnot$
  या प्रकारांसाठी
  \olref[int][sem][rel]{defn:true-at-w} मधील
  $\mSat{M}{\lnot !A}[w]$ ची व्याख्या वापरा;
  म्हणजे $\lnot !A$ हे
  $!A \lif \lfalse$ ने व्याख्यित आहे
  असे मानू नका.
\end{prob}

\begin{prob}
  खालील !!{formula}s अंतःप्रज्ञावादी तर्कशास्त्रात
  !!{derivable} नाहीत, हे दाखवा:
  \begin{enumerate}
    \item $(!A \lif !B) \lor (!B \lif !A)$
    \item $(\lnot\lnot !A \lif !A) \lif (!A \lor \lnot !A)$
    \item $(!A \lif !B \lor !C) \lif \bigl((!A \lif !B)\lor(!A \lif !C)\bigr)$
  \end{enumerate}
\end{prob}

\end{document}
